\section{Tecnologie utilizzate}

\subsection{Backend}

\begin{table}[H]
    \centering
    \begin{tabularx}{\textwidth}{@{}llX@{}}
        \toprule
        \textbf{Tecnologia} & \textbf{Versione} & \textbf{Motivazione della scelta} \\
        \midrule
        Python & 3.10+ &
        Runtime del backend. Linguaggio ad alto livello con un ecosistema
        maturo per lo sviluppo web e la gestione dei dati. \\

        Flask & 3.1 &
        Framework REST leggero e flessibile; non impone strutture rigide,
        permettendo di organizzare il codice secondo il pattern
        \emph{Routes $\to$ Services $\to$ Models}. \\

        Flask-SQLAlchemy & 3.1 &
        ORM che astrae il database in classi Python, garantisce query
        parametrizzate (prevenzione SQL injection) e facilita la
        migrazione futura verso DBMS alternativi. \\

        SQLite & embedded &
        Database relazionale senza server separato; adatto alla fase di
        sviluppo e prototipazione. In produzione è sostituibile con
        PostgreSQL senza modificare le API. \\

        Flask-JWT-Extended & 4.7 &
        Gestione nativa di access token JWT con algoritmo HS256 e
        verifica automatica della firma e della scadenza. \\

        Flask-CORS & 6.0 &
        Abilita le intestazioni CORS necessarie per le richieste
        cross-origin dall'app mobile (e da browser in fase di sviluppo). \\

        werkzeug.security & — &
        Funzioni \code{generate\_password\_hash} e
        \code{check\_password\_hash} basate su PBKDF2-SHA256 con salt
        randomico; inclusa nell'ecosistema Flask, senza dipendenze
        aggiuntive. \\

        \bottomrule
    \end{tabularx}
    \caption{Stack tecnologico del backend.}
\end{table}

\subsection{App Mobile}

\begin{table}[H]
    \centering
    \begin{tabularx}{\textwidth}{@{}llX@{}}
        \toprule
        \textbf{Tecnologia} & \textbf{Versione} & \textbf{Motivazione della scelta} \\
        \midrule
        React Native & 0.81 &
        Framework che genera UI native per iOS e Android da un'unica
        codebase JavaScript/TypeScript, riducendo costi di manutenzione
        rispetto ad app separate. \\

        Expo SDK & 54 &
        Toolchain che semplifica build, hot-reload e accesso alle API
        native (fotocamera, storage sicuro, ecc.) senza configurare
        progetti Xcode o Android Studio. \\

        expo-router & 6 &
        Routing file-based (ogni file in \code{app/} diventa una route),
        con supporto a layout annidati e navigazione type-safe. \\

        TypeScript & — &
        Tipizzazione statica che riduce i bug a runtime e migliora
        l'autocompletamento negli IDE. Particolarmente utile per la
        definizione dei tipi delle risposte REST. \\

        expo-secure-store & — &
        Persiste il JWT nel Keychain (iOS) o Keystore (Android), le
        aree di storage cifrate del sistema operativo; su Web ricade
        su \code{localStorage}. \\

        \bottomrule
    \end{tabularx}
    \caption{Stack tecnologico dell'app mobile.}
\end{table}

\subsection{Protocollo di comunicazione}

Il trasferimento dati tra app mobile e backend avviene tramite
\textbf{HTTP/1.1} con payload \textbf{JSON} codificato UTF-8, in
conformità alla specifica LRCLIB. I client autenticati allegano il token
nell'header \code{Authorization: Bearer <jwt>}; le pubblicazioni
anonime usano l'header custom \code{X-Publish-Token: <prefix>:<nonce>}.

In ambiente di sviluppo le comunicazioni avvengono in HTTP non cifrato
sulla rete locale. In produzione è previsto HTTPS con TLS terminato
da un reverse proxy (nginx o Caddy), come dettagliato nella
sezione~\ref{sec:sicurezza}.
